% Impostazioni comuni: pagina, caratteri, spaziatura e collegamenti.
\usepackage[margin=3cm]{geometry}
\usepackage{fontspec}
\usepackage{lmodern}
\usepackage[italian]{babel}
\usepackage{microtype}
\usepackage{setspace}
\setstretch{1.25}
\usepackage{parskip}
\usepackage{enumitem}
% Elenchi compatti e indivisibili: se non entrano nella pagina passano interi alla successiva.
\makeatletter
\setlist[itemize,enumerate]{nosep, before={\interlinepenalty=10000\relax\@itempenalty=10000\relax}}
\makeatother
\setlength{\emergencystretch}{3em}
% Nessuna riga isolata di un paragrafo in cima o in fondo alla pagina.
% Si impostano a inizio documento perché un pacchetto le riporta a 3000.
\AtBeginDocument{\widowpenalty=10000 \clubpenalty=10000 \displaywidowpenalty=10000}
\setcounter{secnumdepth}{2}
\setcounter{tocdepth}{2}

% Note a piè di pagina sempre in fondo alla pagina, anche quando il testo non la riempie.
\usepackage[bottom]{footmisc}

% Titolo dei capitoli su una riga, con il numero davanti e senza la parola «Capitolo».
\usepackage{titlesec}
\titleformat{\chapter}[hang]{\normalfont\huge\bfseries}{\thechapter}{0.8em}{}
\titlespacing*{\chapter}{0pt}{70pt}{40pt}

\usepackage{graphicx}
\usepackage[table,dvipsnames]{xcolor}

% Diagrammi in TikZ, con gli stessi colori delle tabelle.
\usepackage{tikz}
\usetikzlibrary{arrows.meta,backgrounds,calc,fit,positioning,shapes.geometric}
\definecolor{FiguraTesto}{HTML}{1F2328}
\definecolor{FiguraTenue}{HTML}{59636E}
\definecolor{FiguraBordo}{HTML}{BFBFBF}
\definecolor{FiguraAccento}{HTML}{8C1C2B}
% Seminero di Latin Modern, più stretto del grassetto esteso, per i titoli nei diagrammi.
\newfontfamily\FiguraSeminero{lmromandemi10-regular}[Extension=.otf, ItalicFont=lmromandemi10-oblique]
% Stili comuni: riquadri di raggruppamento con intestazione come nelle tabelle e frecce.
\tikzset{
	partizione/.style={draw=FiguraBordo, fill=TabellaRiga},
	intestazione/.style={draw=FiguraBordo, fill=TabellaIntestazione},
	titolo/.style={font=\footnotesize\bfseries, text=FiguraTesto},
	freccia/.style={draw=FiguraTenue, -{Stealth[length=1.8mm,width=1.4mm]}},
}
% Una voce di diagramma: titolo in seminero e, sotto, un dettaglio più piccolo e grigio.
\newcommand{\voce}[2]{{\fontsize{8.5}{10}\FiguraSeminero #1}\\[2pt]\color{FiguraTenue}#2}
\usepackage{caption}
\captionsetup{font=small}
\usepackage{float}
\floatstyle{plain}
\newfloat{codelisting}{htbp}{lol}[chapter]
\floatname{codelisting}{Codice sorgente}
\usepackage{listings}
\lstset{
	basicstyle=\small\ttfamily,
	breaklines=true,
	columns=fullflexible,
	keepspaces=true,
	showstringspaces=false
}

% Tabelle: stile comune e larghezza complessiva pari al testo.
% Una tabella non si spezza mai tra due pagine: se non entra nello spazio rimasto,
% passa intera alla pagina successiva. Didascalia e identificativo sono argomenti
% dell'ambiente, così la didascalia compare sempre sotto la tabella.
% I margini standard delle colonne allineano sfondi e righe senza patch.
\usepackage{booktabs,array}
\usepackage{needspace}
\definecolor{TabellaIntestazione}{HTML}{E2E5E9}
\definecolor{TabellaRiga}{HTML}{F3F4F6}
\newcolumntype{P}[1]{>{\raggedright\arraybackslash}p{\dimexpr#1\linewidth-2\tabcolsep\relax}}
\newenvironment{tabellatesi}[3]{%
	\def\TabellaDidascalia{#2}%
	\def\TabellaEtichetta{#3}%
	\begin{table}[H]%
	\centering
	\setstretch{1}\fontsize{11}{13.5}\selectfont
	\setlength{\tabcolsep}{8pt}%
	\renewcommand{\arraystretch}{1.3}%
	\setlength{\extrarowheight}{3pt}%
	\setlength{\aboverulesep}{0pt}%
	\setlength{\belowrulesep}{0pt}%
	\arrayrulecolor{black!25}%
	\captionsetup{position=below,skip=6pt}%
	\rowcolors{2}{white}{TabellaRiga}%
	\begin{tabular}{#1}%
}{%
	\end{tabular}%
	\caption{\TabellaDidascalia}\label{\TabellaEtichetta}%
	\end{table}%
}

\usepackage[unicode]{hyperref}
\usepackage{bookmark}
\usepackage{xurl}
\urlstyle{same}
\hypersetup{
	pdftitle={\TitoloTesi},
	pdfauthor={\AutoreTesi},
	pdflang={it-IT},
	colorlinks=true,
	linkcolor=Maroon,
	citecolor=Blue,
	urlcolor=Blue
}
